\section{Conclusiones}

En relación con el objetivo general y la evaluación funcional del software, la investigación demostró que la integración de inteligencia artificial generativa en la creación de código para Godot alcanza una paridad técnica completa con el desarrollo humano tradicional. Las pruebas cuantitativas ejecutadas en la condición ciega mostraron que los usuarios experimentaron niveles equiparables en la experiencia del jugador ($t(29) = 2.33$, $p = 0.0267$). Esto permite concluir que, desde la percepción ciega de los usuarios, el uso de código generado por IA no degrada la experiencia obtenida.

En lo que respecta al análisis del sesgo y la percepción del usuario, el estudio evidenció la existencia de una clara penalización subjetiva condicionada por la información de autoría. Al revelar de manera explícita la presencia de IA en el grupo informado, se registró un decremento acentuado y estadísticamente significativo en la valoración general del videojuego ($t(29) = 8.57$, $p < 0.0001$), respaldado por el contraste de muestras independientes entre grupos ($t(58) = -3.97$, $p = 0.0002$). Este hallazgo confirma que la baja en la calificación de la experiencia del usuario responde a un prejuicio o estigma tecnológico latente hacia las herramientas generativas.

Asimismo, los resultados determinaron que la intensidad de este sesgo está moderada por las preconcepciones y actitudes previas del evaluador respecto a la tecnología. El análisis de correlación de Pearson reflejó una relación inversa moderada ($r = -0.41$, $p = 0.024$) entre la postura crítica o negativa declarada antes de la prueba y la puntuación otorgada en la condición informada. Por lo tanto, se concluye que la resistencia o cautela académica hacia la inteligencia artificial actúa como un factor que predispone al usuario a evaluar el producto con mayor severidad cuando conoce su origen automatizado.

Respecto al contraste de la hipótesis, el valor de $p$ obtenido en la condición informada ($p < 0.0001$), significativamente menor al nivel de significancia establecido ($\alpha = 0.05$), aporta la evidencia empírica suficiente para rechazar la Hipótesis Nula ($H_0$). Se valida formalmente que, si bien el rendimiento del código con IA es equivalente al humano en condiciones ciegas, la valoración final del producto se encuentra significativamente condicionada por el etiquetado y la actitud previa del usuario.

Finalmente, el proyecto cumplió con sus propósitos al evidenciar y cuantificar el rechazo que existe hacia la integración de IA en el desarrollo de videojuegos. Esto mediante una metodología experimental validada que separa la calidad objetiva del software de las percepciones actitudinales del usuario. La suite de prototipos en Godot sirvió como un instrumento efectivo de prueba, demostrando que los criterios de aceptación técnicos se cumplieron en ambas versiones y aportando un marco de referencia empírico para la adopción responsable de la IA en la ingeniería de software interactivo.